\section{Projeto predial de baixa tensão}

\begin{frame}{Projeto predial: sequência correta}
\small
\centering
\begin{tikzpicture}[node distance=0.32cm,every node/.style={font=\scriptsize}]
\tikzset{et/.style={draw=corPrim,fill=corDef!7,rounded corners,minimum width=2.55cm,minimum height=0.72cm,align=center}}
\node[et] (a) {1. Planta e uso\\dos ambientes};
\node[et,right=of a] (b) {2. Pontos e\\cargas};
\node[et,right=of b] (c) {3. Circuitos\\terminais};
\node[et,below=of c] (d) {4. Condutores e\\eletrodutos};
\node[et,left=of d] (e) {5. Proteções\\e aterramento};
\node[et,left=of e] (f) {6. Quadro e\\balanceamento};
\node[et,below=of f] (g) {7. Entrada da\\concessionária};
\node[et,right=of g] (h) {8. Unifilar e\\memorial};
\node[et,right=of h] (i) {9. Verificação\\e as built};
\foreach \u/\v in {a/b,b/c,c/d,d/e,e/f,f/g,g/h,h/i}{\draw[-{Stealth},thick,corPrim] (\u)--(\v);}
\end{tikzpicture}
\vspace{2mm}
\begin{caixaAt}
Não se começa o projeto escolhendo disjuntor, cabo ou padrão de entrada. Esses elementos são consequência da arquitetura, das cargas, dos critérios normativos e das condições reais de instalação.
\end{caixaAt}
\end{frame}

\begin{frame}{Da arquitetura para os pontos elétricos}
\small
\begin{columns}[T,onlytextwidth]
\column{0.50\textwidth}
\begin{caixaDef}[title=Levantamento]
Antes de lançar símbolos elétricos, identificar:
\begin{itemize}
\item dimensões e perímetros dos ambientes;
\item portas, janelas e circulação;
\item bancadas, mobiliário e equipamentos fixos;
\item áreas molhadas, externas e sujeitas a lavagem;
\item posição tecnicamente adequada do QD.
\end{itemize}
\end{caixaDef}
\column{0.48\textwidth}
\begin{caixaProjeto}[title=Depois]
\begin{itemize}
\item calcular os mínimos normativos;
\item acrescentar pontos necessários ao uso real;
\item definir TUE pelas cargas previstas;
\item posicionar interruptores junto aos acessos;
\item lançar rotas compatíveis com a construção.
\end{itemize}
\end{caixaProjeto}
\end{columns}
\end{frame}

\begin{frame}{Iluminação e tomadas: mínimo não é projeto final}
\small
\begin{columns}[T,onlytextwidth]
\column{0.49\textwidth}
\begin{caixaNorma}[title=Previsão mínima de iluminação]
Em habitações, a previsão de carga parte de:
\begin{itemize}
\item 100 VA para os primeiros 6 m$^2$ ou fração;
\item 60 VA para cada acréscimo de 4 m$^2$ inteiros ou em fração.
\end{itemize}
A quantidade real de luminárias decorre do projeto luminotécnico, não dessa potência de previsão.
\end{caixaNorma}
\column{0.49\textwidth}
\begin{caixaNorma}[title=TUG]
A quantidade mínima depende do tipo e do perímetro do ambiente. Cozinha, copa, lavanderia e área de serviço possuem critérios mais exigentes que salas, dormitórios e outros cômodos.
\end{caixaNorma}
\begin{caixaAt}
Não desenhar tomadas apenas para atingir uma contagem: conferir portas, móveis, bancadas, eletrodomésticos e acessibilidade de uso.
\end{caixaAt}
\end{columns}
\end{frame}

\begin{frame}{Divisão de circuitos: critérios defensáveis}
\small
\begin{columns}[T,onlytextwidth]
\column{0.49\textwidth}
\begin{itemize}
\item separar iluminação e tomadas quando tecnicamente adequado;
\item separar cargas de potência elevada ou uso específico;
\item prever circuitos exclusivos para TUE quando requerido;
\item limitar a consequência de uma falha a uma área razoável;
\item facilitar identificação, manutenção e expansão.
\end{itemize}
\column{0.49\textwidth}
\begin{itemize}
\item respeitar ambientes molhados e externos;
\item evitar compartilhamento indevido de neutros;
\item considerar DR/DDR e sua seletividade funcional;
\item distribuir circuitos monofásicos entre as fases;
\item revisar simultaneidade e perfil das cargas.
\end{itemize}
\end{columns}
\begin{caixaObs}
A divisão de circuitos deve aparecer de forma idêntica na planta, no quadro de cargas, no diagrama unifilar e na identificação física do QD.
\end{caixaObs}
\end{frame}

\begin{frame}{Quadro de cargas: não usar números soltos}
\scriptsize
\begin{tabularx}{\textwidth}{c p{2.25cm} c c c c c c X}
\toprule
Circ. & Descrição & Fase & $P$ & $I_b$ & Cabo & DJ & DR & Verificações \\
 & & & W/VA & A & mm$^2$ & A & mA & \\
\midrule
C1 & Iluminação social & C & 720 & 3,27 & 1,5 & 10 & -- & $I_z$, queda, curto \\
C2 & Iluminação íntima & B & 860 & 3,91 & 1,5 & 10 & -- & $I_z$, queda, curto \\
C3 & TUG áreas secas & A & 1100 & 5,00 & 2,5 & 16 & conf. & uso e extensão \\
C4 & TUG cozinha & B & 1900 & 8,64 & 2,5 & 16 & 30 & ambiente molhado \\
C7 & Chuveiro & A & 5500 & 25,0 & 6 & 32 & 30 & dedicado \\
\bottomrule
\end{tabularx}
\vspace{2mm}
\begin{caixaAt}
Os valores acima pertencem ao projeto residencial desenvolvido adiante. Eles não são uma tabela universal para copiar em qualquer residência.
\end{caixaAt}
\end{frame}

\begin{frame}{Planta elétrica: o que precisa existir no desenho}
\small
\begin{columns}[T,onlytextwidth]
\column{0.47\textwidth}
\begin{itemize}
\item paredes, portas, janelas e nomes dos ambientes;
\item cotas ou referência dimensional confiável;
\item luminárias, interruptores e comandos;
\item TUG, TUE, dados e demais sistemas;
\item número do circuito junto aos pontos/rotas;
\item QD em local acessível e coerente.
\end{itemize}
\column{0.51\textwidth}
\begin{itemize}
\item eletrodutos ou rotas com leitura construtiva;
\item seção/quantidade de condutores quando pertinente;
\item alturas e observações de montagem;
\item legenda de símbolos;
\item compatibilização com hidráulica, estrutura, HVAC e dados;
\item detalhes dos locais especiais.
\end{itemize}
\begin{caixaProjeto}
Uma planta elétrica não é um diagrama abstrato colocado sobre retângulos. Ela deve permitir compreender \textbf{onde} instalar, \textbf{o que} instalar e \textbf{a qual circuito} cada ponto pertence.
\end{caixaProjeto}
\end{columns}
\end{frame}

\begin{frame}{Unifilar: fechar a lógica elétrica}
\small
\begin{columns}[T,onlytextwidth]
\column{0.48\textwidth}
\begin{itemize}
\item origem e tensão de fornecimento;
\item medição e proteção de entrada;
\item alimentador e número de condutores;
\item DG e barramentos do QD;
\item DPS e conexão ao sistema de equipotencialização;
\item DR/DDR/RCBO conforme circuitos protegidos.
\end{itemize}
\column{0.50\textwidth}
\begin{itemize}
\item disjuntores terminais;
\item fases atribuídas aos circuitos;
\item neutros corretamente segregados a jusante dos DR;
\item barramento PE contínuo e identificado;
\item seções e correntes nominais;
\item identificação coerente com planta e quadro de cargas.
\end{itemize}
\end{columns}
\begin{caixaAt}
O símbolo de terra isolado não substitui o desenho do PE, do BEP e da equipotencialização. Em projeto didático, esconder essa topologia ensina errado.
\end{caixaAt}
\end{frame}

%=====================================================================
